\section{公民的基本权利和义务}
\begin{enumerate}
  \item 性价比极高。2 - 3 分，朴素的价值观。
  \item 范围封闭，保护开放。区分各个权利主体，一般都是公民，义务主体是国家。
  \item 积极受益 \yjt 积极配合、协助；消极受益 \yjt 国家消极不干预、不侵犯。
  \item 基本权利。主体自然人 \hz 法人。
        \begin{enumerate}
          \item 平等权。
                \begin{enumerate}
                  \item \zd{抽象的平等，不是具体的 \hz 事实意义上的平等}。
                  \item 人格平等。差别对待（\zd{不同主体间}）是否合理判断。年龄、性别、少民华侨。符合尊严、公共利益。
                \end{enumerate}
          \item 政治权利。选举、被选举。
          \item 政治自由。6 个。
                \begin{enumerate}
                  \item 言论。最基础。
                  \item 出版。\zd{批准文号，如书号}。预防追惩相结合。
                  \item 结社。核准登记。民政部和县以上民政部。
                  \item 集会、游行、示威。5 天申请、2 天答复、\zd{公安批}。
                \end{enumerate}
          \item 宗教信仰。无传教、办教自由。
                \begin{enumerate}
                  \item 正常的宗教活动。破坏社会秩序、损害健康、妨害国家的教育制度。
                  \item 宗教团体自主、自办、自传。对外平等交流。
                \end{enumerate}
          \item \zd{人身自由}。狭义，广义。除了\zd{生命权}（内涵），都有明文规定。
          \item 人格尊严。排除任何形式的侵犯（侮辱、诽谤、诬告陷害）。参照民法人格权篇。
          \item 住宅权。居住安全、生活安宁。区分拆迁。
          \item 通信自由、通信秘密。公安和检察，犯罪国安。
          \item 劳动、受教育。既是权利，又是义务。
          \item 获得帮助。弱势群体。
          \item 文化自由。文艺创作、科学研究。
          \item 监督。批评建议 \yjt 任何机关，任何行为；申诉控告检举 \yjt 违法失职，有关机关。
          \item 获得赔偿。
          \item \zd{经社文保障}。
          \item 婚姻家庭保护。
        \end{enumerate}
  \item 基本义务。
\end{enumerate}
